\documentclass[11pt, oneside]{article}   	% use "amsart" instead of "article" for AMSLaTeX format
\usepackage{geometry}                		% See geometry.pdf to learn the layout options. There are lots.
\geometry{letterpaper}                   		% ... or a4paper or a5paper or ... 

\usepackage[parfill]{parskip}    		% Activate to begin paragraphs with an empty line rather than an indent
\usepackage{graphicx}				% Use pdf, png, jpg, or eps§ with pdflatex; use eps in DVI mode
								% TeX will automatically convert eps --> pdf in pdflatex		
\usepackage{amssymb}
%\usepackage{cite}
\usepackage[
backend=biber,
style=verbose,
%style=nature,
sorting=nyt
]{biblatex}
%\addbibresource{../Chemistry.bib}
\bibliography{../Chemistry}

% numbered examples
\usepackage{gb4e}
\usepackage{enumitem}
\usepackage{cancel}
\usepackage{amsmath}

\usepackage{mhchem}
\usepackage{lewis}

\usepackage{url}
\usepackage{hyperref}


% Scientific Laws
\newtheorem{definition}{Definition}
\newtheorem{law}{Law}
\newtheorem{theory}{Theory}
\newtheorem{hint}{Hint}

\title{Module 14: Kinetics}
\author{Mark Peever\\ \texttt{mpeever@gmail.com}}
\date{April 10, 2026}

\begin{document}
\maketitle

\section{Overview}
\begin{enumerate}
\item a \textbf{rate} is a measure of how quickly something changes with respect to time
\item a \textbf{reaction rate} measures how quickly a chemical reaction occurs
\item the \textbf{Rate Equation} allows us to predict a rate for a given reaction
\item chemical reactions \textbf{never actually finish} they just slow to an imperceptible rate
\item \textbf{catalysts} are substances that can affect reaction rates without actually being part of a reaction
\end{enumerate}

\section{Reaction Kinetics}

\begin{definition}[Reaction Rate]\label{def:reaction-rate}
The rate of change of a product in a chemical reaction.
\end{definition}

\begin{definition}[Reaction Rate (reprise)]\label{def:reaction-rate:reprise}
$R = \frac{\Delta [product]}{\Delta t} = - \frac{\Delta [reactant]}{\Delta t} $
\end{definition}

\begin{itemize}
\item \emph{Reaction Rate} measures how quickly the concentration of a chemical reaction's \emph{products} changes with respect to time
\item we'd measure this in $\frac{M}{s}$ or ``molar per second"
\item notice it's the opposite of the rate in change of the concentration of \emph{reactants} (Question: \emph{why}?)
\end{itemize}

\begin{hint}\label{hint:particle-collision}
Chemicals' atoms and/or molecules must collide with each other in order to react.
\end{hint}

\begin{hint}\label{hint:particle-collision-rate}
Higher temperature indicates a substance's particles are moving more rapidly, and therefore collide more frequently.
\end{hint}

\begin{hint}\label{hint:reaction-concentration}
Decreasing concentration decreases chemical reaction rate because there are fewer reactant atoms or molecules to collide with each other.
\end{hint}

\begin{hint}\label{hint:reaction-surface-area}
Increasing surface area of a reactant increases reaction rate because the atoms or molecules of the reactants can more easily mix with each other
\end{hint}

\begin{itemize}
\item chemical reactions generally occur at a greater rate at higher temperatures, because their particles are colliding more frequently at higher temperatures (see Hints \ref{hint:particle-collision} and \ref{hint:particle-collision-rate})
\item chemical reactions generally occur at a greater rate in greater concentrations, because there are more particles to collide (see Hints \ref{hint:reaction-concentration})
\item chemical reactions generally occur at a greater rate in finer substances because they have more surface area (see Hint \ref{hint:reaction-surface-area})
\end{itemize}


\section{The Rate Equation}

\begin{hint}\label{hint:reaction-rate-order}
The \emph{overall order of a chemical reaction} is given by $(x + y)$.
\end{hint}

\begin{hint}\label{hint:reactions-never-finish}
Chemical reactions never actually finish: they just slow to an imperceptible rate.
\end{hint}

\begin{hint}\label{hint:reactions-constant-and-temperature}
The rate of many chemical reactions doubles with every $10 K$ temperature increase.
\end{hint}

\begin{itemize}
\item if we have a chemical reaction like: \ce{a A + b B -> c C + d D} \ldots
\item then we can calculate the rate of the reaction as: $R = k [A]^{x} [B]^{y}$
\footnote{Notice this is a differential equation. We might expect it to see something exponential pop out.}
\item here:
\begin{enumerate}
\item $R$ is the rate of the reaction
\item $k$ is the ``rate constant''
\item $x$ and $y$ are exponents on the concentrations of the reactants
\item $k$, $x$, and $y$ are all experimentally determined
\end{enumerate}
\item notice the reaction rate can vary with things like temperature
\footnote{Notice how this mirrors Gibbs' Free Energy: $\Delta G = \Delta H - T \Delta S$.}
\item in fact, the reaction ``constant" varies with temperature 
\footnote{Which means it's not a constant at all, rather $k = k(T)$.}
(see Hint \ref{hint:reactions-constant-and-temperature})
\item we can add $x$ and $y$ together to get the ``overall order" of a chemical reaction (see Hint \ref{hint:reaction-rate-order})
\item $k$ depends on the chemicals themselves, and is independent of $x$ and $y$
\item we can work out what $k$, $x$, and $y$ are from empirical data
\footnote{But of course, \emph{getting} that data is complicated, see \cite[p. 466]{wile-chem-2}. We'll always work from ``given" data.}
\item if the concentration of a reactant doesn't affect the rate, then we make it's exponent $0$.
\footnote{Because $n^{0} = 1 \forall n \in \mathbb{R} $.}
\item because reactions have the greatest rate when the concentrations of the reactants are the greatest, we expect a reaction to start strong and taper off
\item so reactions tend to have asymptotic behavior (see Hint \ref{hint:reactions-never-finish})
\end{itemize}

\section{Catalysts}

\begin{definition}[Reaction Mechanism]\label{def:reaction-mechanism}
A series of chemical equations that lay out the step-by-step process by which a chemical reaction occurs.
\end{definition}

\begin{definition}[Catalyst]\label{def:catalyst}
A substance that speeds up the rate of a reaction without being used up in the reaction.
\end{definition}

\begin{definition}[Heterogeneous Catalyst]\label{def:catalyst-hetero}
A catalyst in a different phase than the reactants.
\end{definition}

\begin{definition}[Homogeneous Catalyst]\label{def:catalyst-homo}
A catalyst in the same phase as at least one of the reactants.
\end{definition}

\begin{itemize}
\item a \emph{catalyst} is a substance that can be added to a chemical reaction that's not actually a reactant
\item catalysts lower the activation energy of a reaction
\item lowering the activation energy of a reaction increases the rate of the reaction
\item heterogeneous catalysts (see Definition \ref{def:catalyst-hetero}) force reactants together and thus lower activation energy
\item homogeneous catalysts (see Definition \ref{def:catalyst-homo}) participate in a chemical reaction both as a reactant and a product, and alter the reaction mechanism (see Definition \ref{def:reaction-mechanism})
\end{itemize}


\section{Homework}
Review Problems: p. 488 \# 1--10 (not to be turned in)\\
Practice Problems: pp. 489--490 \# 1--10 (due 2026-04-17)\\
 
\nocite{wile-chem-2}
%\bibliography{../Chemistry}
%\bibliographystyle{apalike}

\printbibliography


\section*{Answers to Practice Problems}
Answers to Practice Problems pp. 489--490 \# 1--10:
\begin{enumerate}
\item milk left on the counter has a higher temperature, therefore the reaction rate is higher
\item $R = (7110 \frac{1}{M^{2} s}) [NO]^{2} [O_{2}]$
\item $R = (8.1 \times 10^{-2} \frac{1}{s}) [C_{4} Cl_{9} OH]$
\item $R = 1 \times 10^{-4} \frac{M}{s}$
\item $R = 5.22 \times 10^{-4} \frac{M}{s}$
\item $T_{2} = 65 ^{\circ} C$
\item III
\item heterogeneous, because it is a solid and the reactants are gases
\item 
\begin{enumerate}
\item $3 R$
\item $10 R$
\item no catalyst
\end{enumerate}
\item \ce{Cl} is the \emph{homogeneous} catalyst
\end{enumerate}
\end{document}  